% GENERATED FILE. Edit canonical lessons, then run npm run generate:books.
% canonical-source-hash: fnv1a64:7cb1498e0c10f1fc
% canonical-lessons: SA-C431-arghah, SA-C431-maharghah, SA-C431-samarghah, SA-C431-misrayati, SA-C431-nikasah

\chapter{Dear and Cheap}
\label{ch:sa-dear-and-cheap}
\hlchaptermodality{431}

\begin{chapteropening}
\textbf{By the end of this chapter:} \emph{I can say what is dear, what is cheap and what is real.}

Complete the last lesson of chapter 431: I can say what is dear, what is cheap and what is real.
\end{chapteropening}

\section[\texorpdfstring{\sk{अर्घः} (arghaḥ)}{अर्घः (arghaḥ)}]{\sk{अर्घः} (arghaḥ) — worth; value}

\label{lesson:SA-C431-arghah}



\begin{quote}
Before the new one: say the Sanskrit for a customer, then the Sanskrit for haggles.
\end{quote}

\subsection*{You'll want to know: \sk{अर्घः}}
\textbf{\sk{अर्घः}} — \emph{arghaḥ} — \textquotedblleft{}worth; value\textquotedblright{}.

\begin{quote}
\sk{रत्नस्य} \sk{अर्घः} \sk{महान्} \sk{अस्ति।}
\end{quote}

\begin{quote}
\emph{ratnasya arghaḥ mahān asti.}
\end{quote}

\begin{quote}
The jewel's worth is great.
\end{quote}

A masculine noun.

\begin{cousinweb}
\textbf{\sk{अर्घः}} (\emph{arghaḥ}) comes from the root √arh, \textquotedblleft{}to deserve, to be worth\textquotedblright{}; MW gives \emph{argha} as worth, value, price (MW).
\end{cousinweb}

\begin{grammarlens}[title={What you've built}]
One more word for getting around.
\end{grammarlens}

\subsection*{Your turn}
\begin{itemize}
\raggedright
  \item \emph{Say it:} \emph{arghaḥ}
  \item \emph{Say it:} \emph{arghaḥ}, once more
  \item \emph{Say it:} \emph{ratnasya arghaḥ mahān asti.} — The jewel's worth is great.
  \item \emph{Recall:} say the Sanskrit for a customer, then the Sanskrit for haggles, then say \emph{arghaḥ} again
\end{itemize}

\begin{tcolorbox}[breakable,colback=teal!4,colframe=teal!35!black,title={Before you move on}]
What does \sk{अर्घः} mean? (\textquotedblleft{}Worth; value\textquotedblright{}.) Say it once more.
\end{tcolorbox}
\section[\texorpdfstring{\sk{महार्घः} (mahārghaḥ)}{महार्घः (mahārghaḥ)}]{\sk{महार्घः} (mahārghaḥ) — costly; expensive}

\label{lesson:SA-C431-maharghah}



\begin{quote}
Before the new one: say the Sanskrit for haggles, then the Sanskrit for worth.
\end{quote}

\subsection*{You'll want to know: \sk{महार्घः}}
\textbf{\sk{महार्घः}} — \emph{mahārghaḥ} — \textquotedblleft{}costly; expensive\textquotedblright{}.

\begin{quote}
\sk{सुवर्णं} \sk{महार्घम्} \sk{अस्ति।}
\end{quote}

\begin{quote}
\emph{suvarṇaṃ mahārgham asti.}
\end{quote}

\begin{quote}
Gold is costly.
\end{quote}

\begin{cousinweb}
\textbf{\sk{महार्घः}} (\emph{mahārghaḥ}) is built from \emph{mahā}, \textquotedblleft{}great\textquotedblright{}, and \emph{argha}, \textquotedblleft{}value, price\textquotedblright{}; MW gives the compound as high-priced, costly, very precious (MW).
\end{cousinweb}

\begin{grammarlens}[title={What you've built}]
One more word for getting around.
\end{grammarlens}

\subsection*{Your turn}
\begin{itemize}
\raggedright
  \item \emph{Say it:} \emph{mahārghaḥ}
  \item \emph{Say it:} \emph{mahārghaḥ}, once more
  \item \emph{Say it:} \emph{suvarṇaṃ mahārgham asti.} — Gold is costly.
  \item \emph{Recall:} say the Sanskrit for haggles, then the Sanskrit for worth, then say \emph{mahārghaḥ} again
\end{itemize}

\begin{tcolorbox}[breakable,colback=teal!4,colframe=teal!35!black,title={Before you move on}]
What does \sk{महार्घः} mean? (\textquotedblleft{}Costly; expensive\textquotedblright{}.) Say it once more.
\end{tcolorbox}
\section[\texorpdfstring{\sk{समर्घः} (samarghaḥ)}{समर्घः (samarghaḥ)}]{\sk{समर्घः} (samarghaḥ) — low-priced}

\label{lesson:SA-C431-samarghah}



\begin{quote}
Before the new one: say the Sanskrit for worth, then the Sanskrit for costly.
\end{quote}

\subsection*{You'll want to know: \sk{समर्घः}}
\textbf{\sk{समर्घः}} — \emph{samarghaḥ} — \textquotedblleft{}low-priced\textquotedblright{}.

\begin{quote}
\sk{हट्टे} \sk{शाकं} \sk{समर्घम्} \sk{अस्ति।}
\end{quote}

\begin{quote}
\emph{haṭṭe śākaṃ samargham asti.}
\end{quote}

\begin{quote}
Vegetables are low-priced at the bazaar.
\end{quote}

\begin{cousinweb}
\textbf{\sk{समर्घः}} (\emph{samarghaḥ}) is built from \emph{sam} and \emph{argha}, \textquotedblleft{}price\textquotedblright{}; MW gives \emph{samargha} as cheap (MW).
\end{cousinweb}

\begin{grammarlens}[title={What you've built}]
One more word for getting around.
\end{grammarlens}

\subsection*{Your turn}
\begin{itemize}
\raggedright
  \item \emph{Say it:} \emph{samarghaḥ}
  \item \emph{Say it:} \emph{samarghaḥ}, once more
  \item \emph{Say it:} \emph{haṭṭe śākaṃ samargham asti.} — Vegetables are low-priced at the bazaar.
  \item \emph{Recall:} say the Sanskrit for worth, then the Sanskrit for costly, then say \emph{samarghaḥ} again
\end{itemize}

\begin{tcolorbox}[breakable,colback=teal!4,colframe=teal!35!black,title={Before you move on}]
What does \sk{समर्घः} mean? (\textquotedblleft{}Low-priced\textquotedblright{}.) Say it once more.
\end{tcolorbox}
\section[\texorpdfstring{\sk{मिश्रयति} (miśrayati)}{मिश्रयति (miśrayati)}]{\sk{मिश्रयति} (miśrayati) — mixes; adulterates}

\label{lesson:SA-C431-misrayati}



\begin{quote}
Before the new one: say the Sanskrit for costly, then the Sanskrit for low-priced.
\end{quote}

\subsection*{You'll want to know: \sk{मिश्रयति}}
\textbf{\sk{मिश्रयति}} — \emph{miśrayati} — \textquotedblleft{}mixes; adulterates\textquotedblright{}.

\begin{quote}
\sk{कृपणः} \sk{वणिक्} \sk{दुग्धे} \sk{जलं} \sk{मिश्रयति।}
\end{quote}

\begin{quote}
\emph{kṛpaṇaḥ vaṇik dugdhe jalaṃ miśrayati.}
\end{quote}

\begin{quote}
The stingy merchant mixes water into the milk.
\end{quote}

Past (he …): \sk{सः} \sk{दुग्धे} \sk{जलं} \sk{मिश्रितवान्}, \emph{saḥ dugdhe jalaṃ miśritavān}.

\begin{cousinweb}
\textbf{\sk{मिश्रयति}} (\emph{miśrayati}) is built on \emph{miśra}, \textquotedblleft{}mixed\textquotedblright{}, as MW notes; MW gives \emph{miśrayati} as to mix or blend, with the participle \emph{miśrita}; MW compares Latin \emph{miscere} and English \emph{mix} (MW).
\end{cousinweb}

\begin{grammarlens}[title={What you've built}]
One more word for getting around.
\end{grammarlens}

\subsection*{Your turn}
\begin{itemize}
\raggedright
  \item \emph{Say it:} \emph{miśrayati}
  \item \emph{Say it:} \emph{miśrayati}, once more
  \item \emph{Say it:} \emph{kṛpaṇaḥ vaṇik dugdhe jalaṃ miśrayati.} — The stingy merchant mixes water into the milk.
  \item \emph{Recall:} say the Sanskrit for costly, then the Sanskrit for low-priced, then say \emph{miśrayati} again
\end{itemize}

\begin{tcolorbox}[breakable,colback=teal!4,colframe=teal!35!black,title={Before you move on}]
What does \sk{मिश्रयति} mean? (\textquotedblleft{}Mixes; adulterates\textquotedblright{}.) Say it once more.
\end{tcolorbox}
\section[\texorpdfstring{\sk{निकषः} (nikaṣaḥ)}{निकषः (nikaṣaḥ)}]{\sk{निकषः} (nikaṣaḥ) — a touchstone}

\label{lesson:SA-C431-nikasah}



\begin{quote}
Before the new one: say the Sanskrit for low-priced, then the Sanskrit for mixes.
\end{quote}

\subsection*{You'll want to know: \sk{निकषः}}
\textbf{\sk{निकषः}} — \emph{nikaṣaḥ} — \textquotedblleft{}a touchstone\textquotedblright{}.

\begin{quote}
\sk{स्वर्णकारः} \sk{निकषे} \sk{सुवर्णं} \sk{घर्षति।}
\end{quote}

\begin{quote}
\emph{svarṇakāraḥ nikaṣe suvarṇaṃ gharṣati.}
\end{quote}

\begin{quote}
The goldsmith rubs the gold on the touchstone.
\end{quote}

A masculine noun.

\begin{cousinweb}
\textbf{\sk{निकषः}} (\emph{nikaṣaḥ}) is built from \emph{ni} and the root √kaṣ, \textquotedblleft{}to rub\textquotedblright{}; MW gives \emph{nikaṣa} as rubbing in and as the touchstone, the stone on which gold is rubbed to test it (MW).
\end{cousinweb}

\begin{grammarlens}[title={What you've built}]
One more word for getting around.
\end{grammarlens}

\subsection*{Your turn}
\begin{itemize}
\raggedright
  \item \emph{Say it:} \emph{nikaṣaḥ}
  \item \emph{Say it:} \emph{nikaṣaḥ}, once more
  \item \emph{Say it:} \emph{svarṇakāraḥ nikaṣe suvarṇaṃ gharṣati.} — The goldsmith rubs the gold on the touchstone.
  \item \emph{Recall:} say the Sanskrit for low-priced, then the Sanskrit for mixes, then say \emph{nikaṣaḥ} again
\end{itemize}

\begin{tcolorbox}[breakable,colback=teal!4,colframe=teal!35!black,title={Before you move on}]
What does \sk{निकषः} mean? (\textquotedblleft{}A touchstone\textquotedblright{}.) Say it once more.
\end{tcolorbox}
